\section{Introduction}

Probabilistic transition systems can be represented by different kinds of models. Markov chains, for example, are a familiar model in this setting, while probabilistic automata provide a broader framework (see, e.g.,~\cite{Segala95}). Logical languages can then be used to describe system properties and reason about system behavior. These include temporal logics, used in PRISM to specify properties of Markov chains~\cite{Kwiatkowska2007} and to reason about knowledge~\cite{HT93}. Finally, Probabilistic Modal Logics are specifically tailored to  reasoning about probabilistic Kripke structures (see, e.g.,~\cite{ALDINI2024109167}) and provide a more general view on transition systems.

Relation-changing formalisms provide another perspective on transition systems. They are discrete-state models in which the accessibility relation can change in response to an event, such as traversing an edge. These changes can be semantically encoded in the model or introduced syntactically through operators in the language. In switch graphs, for example, changes to the relation are encoded in the model structure itself~\cite{GabbayM12}. Classical modal logic can be used to explore such a model without itself changing the relation. Instead, the prescribed updates occur when designated edges are traversed. Alternatively, one can start with a standard Kripke structure and use a logic whose syntax includes relation-changing operators. Sabotage and bridge logics, for example, include operators that modify the model during formula evaluation, so that the remaining subformula is evaluated on the modified model~\cite{ArecesFH15,Aucher&vanBenthem&Grossi18,Fervari14}.
This kind of dynamic behavior can be also found in many applications of logics interpreted over graph-like structures, such as logics of public announcements and its relatives in epistemic logic (see, e.g.,~\cite{Plaza89:lopc,FervariV19}), logics for resource allocation such as separation logics (see, e.g.,~\cite{OHearn19,DemriD15})  and epistemic planning~\cite{BolanderA11}, to name a few.
% \textcolor{red}{Should we mention logics with public announcements?}

In recent years, the relation-changing paradigm has been extended to several classes of models, including paraconsistent~\cite{CostaFM23}, weighted~\cite{FRMC19}, fuzzy~\cite{FFM25}, and labeled models~\cite{CFMS25,FMB18}. Applications and computational tools have also emerged. For instance, \cite{FRMC19} presents a tool for a probabilistic setting based on a particular choice of weights. In \cite{MFAPCB24}, a related approach is applied in a practical setting, illustrating its potential for real-world applications. Other case studies include~\cite{FB18,SMF21}.

This work advances this line of research by introducing relation-changing operators for probabilistic systems, rather than restricting attention to predefined models with limited expressive power. Although relation-changing approaches have also been studied in fuzzy settings~\cite{FFM25}, the probabilistic setting is not simply a special case of the fuzzy one. In particular, the outgoing transition probabilities at each state must sum to one, so changing one transition may require corresponding adjustments to others. This normalization constraint also poses new challenges for the metatheoretic analysis of the resulting logics. 

In this article we present a probabilistic setting for relation-changing logics, and study associated probabilistic bisimulation notions. We expand the results presented in~\cite{FFM25b}, by providing full proofs, more detailed discussions, examples of use, and incoporating new results. In particular, one of the novel contributions herein with respect to~\cite{FFM25b} is the proof of Hennessy-Milner results for the operators we studied.

\paragraph{Outline} In Section~\ref{sec:pml} we introduce the syntax, semantics and bisimulation notion for the basic probabilistic modal logic that we consider the basis of our framework. Section~\ref{sec:prcml} is devoted to exploring probabilistic versions of sabotage and bridge modalities. Then, in Section~\ref{sec:bisimulation}, we present their associated notions of bisimulations and show their corresponding invariance theorems. Therein, we also investigate respective Hennessy-Milner classes and prove their adequacy. We conclude in Section~\ref{sec:final} with some remarks and future lines of work.